\notechapter{N-20220315}{Quadratic Parameters, Tangency, and Constrained Extrema}

\section{A scattered algebraic note}

The original page is titled ``some scattered things.''  The common theme is that elementary-looking algebraic problems often hide a parameter, a tangency condition, or an extremum under constraint.  I rewrite the problems in a more systematic form.

\section{A quadratic with parameters}

Let
\[
  f(x)=ax^2+bx+c,\qquad a<b<c,
\]
and suppose the graph passes through \(A=(1,0)\) and \(B=(m,-a)\).  The condition \(f(1)=0\) gives
\[
  a+b+c=0.
\]
The condition \(f(m)=-a\) is equivalent to saying that \(f(x)+a\) has \(x=m\) as a root:
\[
  am^2+bm+c+a=0.
\]
Using \(c=-a-b\), this becomes
\[
  am^2+bm-b=0,\qquad am^2=-b(m-1).
\]
The axis of symmetry is
\[
  x_0=-\frac{b}{2a}.
\]
Thus the parameter \(m\) and the axis \(x_0\) are tied by
\[
  x_0=\frac{m^2}{2m-2}.
\]
Equivalently,
\[
  m^2=(2m-2)x_0,\qquad m^2-2x_0m+2x_0=0.
\]
Solving for \(m\) gives
\[
  m=x_0\pm\sqrt{x_0^2-2x_0}.
\]
The original calculation uses this relation to examine endpoints and stationary points of \(m\) as a function of \(x_0\).  The conclusion is that the extremal parameter occurs at
\[
  x_0=0\quad\text{or}\quad x_0=-\frac12,
\]
and the relevant value is
\[
  m=\frac{\sqrt5-1}{2}.
\]
This is a small example of a recurring trick: convert the algebraic parameter into a geometric parameter, usually the axis, vertex, or tangency point.

\section{A parabola and a hyperbola}

The second problem considers the parabola
\[
  y=x^2-2ax+a^2-4=(x-a)^2-4
\]
and the rectangular hyperbola
\[
  y=\frac{k}{x},\qquad k>0,
\]
in the first quadrant.  The intersection equation is
\[
  x^2-2ax+a^2-4=\frac{k}{x},\qquad x>0,
\]
or
\[
  x^3-2ax^2+(a^2-4)x-k=0.
\]
Questions about ``only one intersection'' are therefore tangency or root-counting questions for this cubic, but the geometry is usually cleaner: one restricts to the right side of the parabola, because the branch \(k/x\) lies in the first quadrant.

\begin{remark}
The hidden principle is that a graph-intersection problem may become either a polynomial root-counting problem or a tangency problem.  The latter is often more geometric and less computational.
\end{remark}

\section{A constrained extremum}

The third problem asks for extrema of
\[
  x+y
\]
under the constraint
\[
  x^2+y^3=2.
\]
Using Lagrange multipliers for
\[
  F(x,y,\lambda)=x+y+\lambda(x^2+y^3-2),
\]
we get
\[
  1+2\lambda x=0,\qquad
  1+3\lambda y^2=0,\qquad
  x^2+y^3=2.
\]
Eliminating \(\lambda\) gives
\[
  2x=3y^2.
\]
Hence
\[
  x=\frac32y^2.
\]
Substitution into the constraint gives
\[
  \frac94y^4+y^3=2,
\]
so the stationary points come from the quartic
\[
  9y^4+4y^3-8=0.
\]
The numerical plot in the original note marks approximately
\[
  A=(1.15195,2.02829),\qquad B=(1.82947,0.72509),
\]
for the function
\[
  h(x)=x+(2-x^2)^{1/3}.
\]
These are the local maximum and local minimum along the real branch.

\begin{figure}[H]
\centering
\begin{tikzpicture}[scale=1.2,>=Stealth,every node/.style={fill=white,inner sep=1.2pt}]
  \draw[->] (-1.5,0)--(4.3,0) node[right] {$x$};
  \draw[->] (0,-1.25)--(0,3.2) node[above] {$y$};
  \draw[smooth,thick] plot coordinates {(-1.2,-0.45) (-0.5,0.55) (0,1.26) (0.6,1.78) (1.15,2.03) (1.45,1.55) (1.83,0.73) (2.5,0.95) (4.0,1.65)};
  \fill (1.15,2.03) circle (0.04);
  \node[above left=2pt] at (1.15,2.03) {$A\approx(1.15,2.03)$};
  \fill (1.83,0.73) circle (0.04);
  \node[below right=2pt] at (1.83,0.73) {$B\approx(1.83,0.73)$};
  \node[below] at (2.25,-0.55) {$y=x+(2-x^2)^{1/3}$};
\end{tikzpicture}
\caption{Reconstruction of the graph used for the constrained extremum problem.}
\end{figure}


\section{Tangency as a double-root condition}

A parameter problem involving a quadratic often asks when a line is tangent to a parabola or when an equation has exactly one solution.  Algebraically, this is a discriminant condition:
\[
  \Delta=b^2-4ac=0.
\]
Geometrically, the line touches the parabola at one point.  The same idea appears in constrained extrema: a level curve just touches a constraint curve at an optimum.

\section{Completing the square}

For a quadratic
\[
  ax^2+bx+c,\qquad a\ne0,
\]
completing the square gives
\[
  a\left(x+\frac b{2a}\right)^2+c-\frac{b^2}{4a}.
\]
This form reveals the vertex, minimum or maximum value, and symmetry axis.  It is usually more informative than the expanded form.

\section{Expanded synthesis and advanced context}

The three examples are connected by the same idea: reduce a complicated statement to a structural condition.  The axis of a quadratic, the tangency of two graphs, and the Lagrange multiplier equations are all ways of saying that an extremum occurs when a first-order freedom disappears.


\paragraph{Expanded synthesis.}
For this chapter, the elementary material should be read as a study of what remains invariant when coordinates change. The earlier sections (`A scattered algebraic note`, `A quadratic with parameters`, `A parabola and a hyperbola`, `A constrained extremum`, `Tangency as a double-root condition`) compute with arrays, bases, pivots, determinants, or eigenvectors; the deeper object is a linear map together with the spaces it acts on. A matrix is a representation of that map after bases have been chosen. This is why formulas such as
\[
  A' = P^{-1}AP,\qquad \widetilde A = T^{-1}AS
\]
are not cosmetic: they distinguish an intrinsic morphism from a coordinate description.

The structural hierarchy is as follows. Kernels record what a map forgets, images record what it reaches, cokernels record obstructions to solvability, and quotients record information after an equivalence has been imposed. Determinants act on the top exterior power and therefore measure oriented volume; traces are invariant under cyclic rearrangement and become characters in representation theory; adjoints depend on an inner product and lead to self-adjoint spectral theory.

A useful way to return to the computations is to ask, for every formula, which invariant it is measuring. Row reduction measures rank and kernel; diagonalization measures decomposition into invariant subspaces; Gram--Schmidt measures orthogonal projection; positive definiteness measures ellipsoid geometry; tensor notation measures how components transform. The elementary methods are therefore the finite-dimensional laboratory in which duality, quotienting, functoriality, and spectral decomposition can be seen clearly.


\section{Pedagogical repair: what these exercises were really training}

A useful way to reread this note is not as three unrelated problems, but as three appearances of the same principle: an extremum is usually where one degree of freedom has been used up.  In the quadratic problem, the hidden freedom is the position of the axis.  In the parabola--hyperbola problem, the hidden freedom is whether the two graphs cross twice or merely touch.  In the constrained problem \(x^2+y^3=2\), the hidden freedom is the tangent direction along the curve.

For the constrained extremum, one should not memorize Lagrange multipliers as a black box.  The geometric sentence is: at an extremum of \(f\) restricted to \(g=0\), the level curve of \(f\) is tangent to the constraint curve.  Algebraically this is
\[
  \nabla f=\lambda\nabla g.
\]
This explains why the multiplier equations feel like ``parallel gradient'' equations.  The formula is not a trick; it is the analytic translation of tangency.
